\documentclass[12pt,a4paper]{article}
\usepackage{graphicx}
\usepackage{amsmath}
\usepackage{hyperref}
\usepackage{geometry}
\geometry{margin=2.5cm}
\emergencystretch=3em

\title{\textbf{CDFD Runtime Paper 06: CLI-First Platform and Web Architecture}}
\author{Steve Bico Mujjabi, MD\\
\small Independent Researcher\\
\small Founder, Vura Labs\\
\small Kampala, Uganda\\
\small ORCID: \href{https://orcid.org/0009-0001-0556-5516}{0009-0001-0556-5516}}
\date{June 2026}

\begin{document}
\maketitle

\begin{abstract}
\noindent \textbf{Executive synthesis.} The platform direction is CLI first,
web app second. The CLI is for serious engine work: reproducible simulations,
paper diagnostics, CDFL language workbench commands, debugging, automation, and
saved results. An optional web studio (\texttt{webapp/}) is a visual layer over
the same backend, not a separate physics system. Optional LLM provider calls are
interpretive tools above saved deterministic outputs. This paper defines those
boundaries and explains how CLI outputs flow into optional visual interfaces.
\end{abstract}

\paragraph{Greek-symbol convention.}
Unless a manuscript states a tighter equation-local meaning, Greek symbols are local CDFD variables or coefficients, not universal constants. Here $\Phi$ denotes driving flux, $\Psi_s$ denotes the adaptive operating ratio, $\Omega$ denotes a domain or state space, and $\Lambda$ denotes the Life Number where used. The coefficients $\alpha$, $\beta$, and $\gamma$ denote accumulation or reinforcement, relaxation or decay, and diffusion, coupling, or redistribution. The symbols $\delta$ and $\epsilon$ denote perturbation or tolerance scales, while $\eta$, $\lambda$, $\mu$, $\nu$, $\rho$, $\sigma$, $\tau$, and $\phi$ denote local efficiency, rate, baseline, frequency, density or correlation, noise or variance, time-scale, and phase or field terms. Subscripts restrict any coefficient to the named pathway, tissue, domain, or experiment.


\begin{figure}[ht]
\centering
\fbox{\begin{minipage}{0.92\textwidth}
\centering
\textbf{Runtime evidence path}\\[0.4em]
Prior CDFD mathematics $\longrightarrow$ CDFL/runtime state $\longrightarrow$ bounded output $\longrightarrow$ falsification target.
\end{minipage}}
\caption{Conceptual schematic for the runtime papers. The engine translates the mathematical frame from the prior CDFD papers into executable CDFL/runtime diagnostics; candidate outputs remain tied to explicit tests before any stronger claim is made.}
\end{figure}

\section{Prior CDFD Basis}
This manuscript does not rederive the mathematical core of CDFD. It uses the
Mujjabi Adaptive Operating Ratio and the constrained-vacuum notation introduced
in Part I \cite{MujjabiI2026}, the torus-knot hierarchy and boundary-mode work
from the Part I sequence \cite{MujjabiVI2026,MujjabiIX2026}, the public balance
law developed in the Part I capstone \cite{MujjabiX2026}, the Part I transport
tests \cite{MujjabiXII2026}, and the Life Number and tri-regime biology bridge
from Part II \cite{MujjabiPartII2026}. The runtime papers therefore act as an
execution layer: they keep the mathematics connected to commands, outputs,
falsification tests, and domain-specific limits.


\section{Ordering Principle}
The platform order is:
\begin{enumerate}
    \item core engine;
    \item CLI;
    \item API or service wrappers;
    \item web app.
\end{enumerate}
The current repository has the engine and CLI as primary surfaces. The Streamlit
studio under \texttt{webapp/} (CDFD Runtime Studio) is optional: it is a visual
layer over \texttt{runtime/runner.py}, the same kernel,
\texttt{runtime/diagnostics.py}, CDFL tooling, and domain adapters, and it does
not duplicate physics logic.

\section{CLI Responsibilities}
The CLI supports:
\begin{itemize}
    \item listing registered domains (\texttt{cdfd.py domains});
    \item Part II diagnostics export (\texttt{cdfd.py diagnostics});
    \item running domain demos with JSON payloads or OOL source scenarios;
    \item validating, executing, linting, formatting, and inspecting CDFL models
    through \texttt{cdfd.py cdfl ...};
    \item exporting JSON result envelopes;
    \item doctor, gallery, compare, report, and explain commands for public
    smoke tests and reproducible run artifacts;
    \item optional provider inventory, status, dry-run prompt audit, and
    research interpretation through \texttt{cdfd.py llm ...};
    \item preserving finite-value and provenance metadata on every command.
\end{itemize}
This is the credibility layer. It gives papers exact commands instead of vague
claims that "the engine was run."

\section{CDFD Runtime Studio}
The public Streamlit studio (\texttt{python -m webapp.run\_server}) already
provides an initial visual environment: regime-banded trajectories, field
heatmaps, Part II source-mix panels, a Life Number calculator, Granger-style
causal graphs from run history, the \texttt{origins\_of\_life} domain explorer,
and a CDFL Workbench. The workbench lets a user edit a CDFL model, validate it,
lint it, run it, format it, inspect its AST, and export JSON/Markdown/HTML
result envelopes through the same backend used by the CLI. Future CDFD Studio
releases may add richer model building, coherence monitoring, and deeper result
browsing on top of the same runtime backend.

The first public demos remain non-medical and visual:
\begin{itemize}
    \item heat diffusion and thermal flow;
    \item traffic or infrastructure routing;
    \item resource-flow networks;
    \item predator-prey/ecology models;
    \item photonic or adaptive-surface demos.
\end{itemize}
Medical applications require stricter validation, privacy, and user-safety
boundaries and are not the first public web surface.

\section{API Boundary}
Any API or service wrapper exposes command-like operations:
\begin{itemize}
    \item validate a CDFL model;
    \item run a CDFL model;
    \item lint, format, and inspect CDFL AST summaries;
    \item retrieve saved result envelopes;
    \item inspect finite audits;
    \item list experiments and selected outputs;
    \item export reports.
\end{itemize}
Each endpoint calls \texttt{runtime/runner.py} or a later service wrapper around
it. The service layer may add authentication, queues, and storage, but not a
second physics implementation.

\section{API Keys and the AI Boundary}
Provider keys belong at the application/provider boundary, not inside the
physics kernel. The runtime exposes optional \texttt{cdfd.py llm} commands
above saved deterministic outputs: provider inventory, provider-key status,
dry-run prompt audit, and provider-backed research interpretation. Supported
provider shapes are explicit rather than universal. Raw keys are never printed,
and \texttt{cdfd.py auth} remains only a compatibility alias for provider-key
status.

This split keeps the engine clean: \texttt{engine/}, CDFL execution, and domain
adapters produce deterministic result envelopes; \texttt{runtime/llm.py}, the
CLI, web layer, VOS, or other governed products may interpret saved outputs
through provider adapters. LLM commentary is therefore a separate interpretive
artifact, not an input to the engine equations.

\section{Vacuum OS}
Vacuum OS (VOS) is the proposed operating system layer above CDFD Runtime. VOS
is not a second engine. It is the orchestration shell for projects that need
apps, CDFL packages, user spaces, experiment queues, provenance stores, and
controlled AI assistance. In this architecture, CDFD Runtime remains the
mathematical kernel; VOS becomes the operating environment that launches models,
manages keys, schedules runs, stores outputs, and connects approved applications
to the same CDFD command backend.

\section{Experiment Visualization}
The experiment outputs already define future panels:
\begin{itemize}
    \item vacuum hysteresis: $M_s$ residual heatmap and recovery curve;
    \item knot stability: stability-index trace toward the $\chi^*$ attractor;
    \item OOL phase: $\Lambda=1$ phase boundary;
    \item frontier sweep: pattern distribution and candidate triage table;
    \item gigamarathon sweep: domain coverage dashboard.
\end{itemize}
The public Streamlit dashboard under \texttt{webapp/dashboard.py} implements an
initial subset: regime-banded $\Psi_s$ trajectories, $\Phi$ and $C$ field
heatmaps, Part II aromatic source-mix bars, a Life Number calculator,
Granger-style causal edges from run history, the \texttt{origins\_of\_life}
domain explorer, and the CDFL Workbench. The web app presents these as model
diagnostics and falsification targets, not as empirical proof.

\section{Conclusion}
The platform makes one thing unmistakable: CDFD is executable. The CLI
establishes that discipline. The optional web studio makes the same discipline
visible without weakening the runtime boundary or replacing command-line
reproducibility.
\begin{thebibliography}{9}
\bibitem{MujjabiI2026}
Steve Bico Mujjabi, MD. \emph{Vortex Stability in a Constrained Transport Vacuum and the Origin of the Fine-Structure Constant}. CDFD Part I Paper I, preprint, 2026.

\bibitem{MujjabiVI2026}
Steve Bico Mujjabi, MD. \emph{The Universal Torus Knot Hierarchy: $Z_n$ Symmetry, the Cinquefoil Family, and the General Structure of CDFT Particle Families}. CDFD Part I Paper VI, preprint, 2026.

\bibitem{MujjabiIX2026}
Steve Bico Mujjabi, MD. \emph{Even-$n$ Torus Knots in CDFT: The Exclusion of $T(2,2)$, the Extension to $T(2,2m)$, and the Anti-Phase Mass Scaling Law}. CDFD Part I Paper IX, preprint, 2026.

\bibitem{MujjabiX2026}
Steve Bico Mujjabi, MD. \emph{The Vacuum Equation of State: Deriving Torus Knot Self-Consistency from the Public CDFD Balance Law $\Psi = \Phi/C$}. CDFD Part I Paper X, preprint, 2026.

\bibitem{MujjabiXII2026}
Steve Bico Mujjabi, MD. \emph{Friction, Superconductivity, Plasma States, and Transport Mysteries: Observable Tests of Adaptive Constraint}. CDFD Part I Paper XII, preprint, 2026.

\bibitem{MujjabiPartII2026}
Steve Bico Mujjabi, MD. \emph{CDFD Part II: Origins of Life and Tri-Regime Bioenergetics}. Public release archive, 2026, \url{https://doi.org/10.5281/zenodo.20264779}.
\end{thebibliography}
\end{document}
